% TOP_PROBLEM: {"problemId": "problem.keating-snaith-conjecture-on-moments-of-the-riemann-zeta-function", "canonicalTitle": "Keating-Snaith Conjecture on Moments of the Riemann Zeta Function", "releaseRank": 162, "primaryDomain": "number_theory_arithmetic_geometry", "primaryDomainLabel": "Number theory & arithmetic geometry", "displayStatus": "open_with_solved_subcases", "releaseStatus": "open_with_solved_subcases"}
% !TeX program = lualatex
% Definition and sources only; this file is not an LLM solution attempt.
\documentclass[11pt]{article}
\usepackage[margin=25mm]{geometry}
\usepackage{fontspec}
\setmainfont{Latin Modern Roman}
\usepackage{amsmath,amssymb,mathtools}
\usepackage{unicode-math}
\setmathfont{Latin Modern Math}
\usepackage{xurl,hyperref}
\hypersetup{colorlinks=true,urlcolor=blue}
\setcounter{secnumdepth}{0}
\setlength{\parindent}{0pt}
\setlength{\parskip}{5pt}
\setlength{\emergencystretch}{3em}
\begin{document}
\section*{\texorpdfstring{Keating-Snaith Conjecture on Moments of the Riemann Zeta Function}{Keating-Snaith Conjecture on Moments of the Riemann Zeta Function}}\label{162}
\subsection{Definitions and mathematical statement}
For fixed real $k>0$, define $I_k(T)=\int_0^T|\zeta(\tfrac12+it)|^{2k}{\,\mathrm{d}} t$. Let $G$ be the standard Barnes $G$-function, normalized by $G(1)=1$ and its usual canonical product, with $G(z+1)=\Gamma(z)G(z)$. Define
\[
 a_k=\prod_{p\ \mathrm{prime}}(1-p^{-1})^{k^2}
 \sum_{m=0}^{\infty}\left(\frac{\Gamma(m+k)}{\Gamma(k)m!}\right)^2p^{-m}.
\]
The assertion is
$I_k(T)\sim a_kG(1+k)^2G(1+2k)^{-1}T(\log T)^{k^2}$ as $T\to\infty$. The parameter $k$ is fixed before this limit, and the target includes the leading constant, not just the logarithmic exponent. The functional equation alone is not being used as a unique definition of $G$.
\par\smallskip\textit{Formulation and concept references:} \href{https://empslocal.ex.ac.uk/people/staff/mrwatkin/zeta/keating-snaith.pdf}{[S1]}.

\subsection{Short English statement}
Do all positive real moments of the zeta function on its central line have the precise leading asymptotic predicted by the arithmetic product and random-matrix factor?
\subsection{Sources}
\begin{itemize}
\item[S1]J. P. Keating and N. C. Snaith, "Random Matrix Theory and zeta(1/2+it)," Communications in Mathematical Physics 214 (2000), 57-89.
\url{https://empslocal.ex.ac.uk/people/staff/mrwatkin/zeta/keating-snaith.pdf}.
\textit{Formulation source.}
\end{itemize}
\end{document}
